% =====================================================
% ROTEIRO PRÁTICO 05
% =====================================================

\section{Roteiro Prático 05 — Retomada de Atividades e Diagnóstico ao Vivo}

% =====================================================
% OBJETIVOS
% =====================================================

\begin{boxObjetivos}
\textbf{Objetivos da Aula}

Ao final desta aula, cada grupo de projeto será capaz de:

\begin{itemize}
\item Conhecer a Profª. Karina e compreender as diretrizes do 2º semestre;
\item Visualizar o diagnóstico real das entregas da turma a partir da apresentação em slides;
\item Zerar na planilha oficial todas as lacunas críticas apontadas no diagnóstico (Pitch, Diagrama, Experimentos);
\item Classificar o estágio de maturidade do projeto na coluna \textbf{STATUS};
\item Apresentar em plenária a situação atual da sua linha e definir as 3 tarefas prioritárias da Sprint 1.
\end{itemize}
\end{boxObjetivos}

% =====================================================
% INTRODUÇÃO & RECURSOS
% =====================================================

\subsection{Introdução}

Bem-vindos de volta ao 2º semestre do Projeto Integrador II! 

Nesta aula prática de retomada, faremos um alinhamento direto com a equipe docente (\textbf{Profª. Karina e Prof. André Moraes}) e utilizaremos a \textbf{Planilha Oficial de Acompanhamento} como nosso instrumento central de trabalho ao vivo.

\vspace{0.3cm}

\RepositorioQRCode{https://docs.google.com/spreadsheets/d/1w7wAwPBdNTfh7lvynoP7x35S__lvRyuTjRqP26oEhNQ}{https://docs.google.com/spreadsheets/d/1w7wAwPBdNTfh7lvynoP7x35S__lvRyuTjRqP26oEhNQ}

% =====================================================
% PARTE 1
% =====================================================

\subsection{Parte 1 — Boas-Vindas e Projeção dos Slides de Diagnóstico (15 minutos)}

\begin{boxAtividade}[Atividade 1 — Apresentação da Profª. Karina e Slides da Turma]

Abertura da aula com projeção do material de slides web (\texttt{retomada-de-atividades-status-projetos/index.html}):

\begin{enumerate}
\item \textbf{Apresentação Docente:} Recepção à Profª. Karina e visão geral da dinâmica do 2º semestre;
\item \textbf{Exibição do Raio-X da Turma:} Projeção dos slides contendo os indicadores de entregas dos 10 projetos da turma INF-2024;
\item \textbf{Revisão do Guia de Colunas (Slide 5):} Esclarecimento sobre o significado e a importância de cada entregável da planilha (Objetivo, GitHub, Overleaf, Diagrama, Pitch, Exp 1-4 e Bloco STATUS);
\item \textbf{Identificação dos Gargalos:}
  \begin{itemize}
      \item \textbf{Pitch Vídeo (3 min):} 90\% da turma pendente (apenas o grupo SMC cadastrou);
      \item \textbf{Diagrama Canva e Experimento 1:} Pendência crítica nos projetos \textit{JifTV} e \textit{RasPorta};
      \item \textbf{Agendamento Externo:} Confirmação da reunião marcada do grupo \textit{Curso Popular de Informática} com a Profª. Karina.
  \end{itemize}
\end{enumerate}

\end{boxAtividade}

% =====================================================
% PARTE 2
% =====================================================

\subsection{Parte 2 — Tarefa Prática 1: Zerar Pendências na Planilha ao Vivo (35 minutos)}

\begin{boxAtividade}[Atividade 2 — Preenchimento e Regularização da Linha do Grupo]

Cada grupo deverá abrir a \textbf{Planilha Oficial de Acompanhamento} e executar a regularização imediata da sua linha:

\begin{enumerate}
\item \textbf{Cadastrar Link do Pitch em Vídeo (3 min):} Todos os 9 grupos pendentes devem colar o link do vídeo de apresentação da solução na coluna \textbf{Pitch apresentação};
\item \textbf{Regularizar Diagrama e Experimento 1:} Os grupos \textit{JifTV} e \textit{RasPorta} devem preencher o link do \textbf{Diagrama Canva} e registrar a descrição sumária do \textbf{Experimento 1};
\item \textbf{Validar Link do Relatório Overleaf:} Verificar se a coluna \textbf{URL Relatório} contém o link funcional do artigo científico com permissão de edição para os professores;
\item \textbf{Atualizar o Bloco STATUS Vivo:} Preencher diretamente nas colunas correspondentes:
  \begin{itemize}
      \item \textbf{AVANÇOS:} O que já foi desenvolvido e funciona atualmente;
      \item \textbf{DIFICULDADES:} Gargalos técnicos, de infraestrutura ou de equipe;
      \item \textbf{TECNOLOGIAS:} As tecnologias e linguagens efetivamente utilizadas.
  \end{itemize}
\item \textbf{Classificar a Coluna STATUS:} Selecionar obrigatoriamente uma das opções padronizadas:
\begin{center}
\textbf{\small [ Diagnóstico \ \textbar \ Planejamento \ \textbar \ Execução \ \textbar \ Experimentos \ \textbar \ Consolidação \ \textbar \ Finalizado ]}
\end{center}
\end{enumerate}

\end{boxAtividade}

% =====================================================
% PARTE 3
% =====================================================

\subsection{Parte 3 — Tarefa Prática 2: Apresentação em Plenária Linha a Linha (30 minutos)}

\begin{boxAtividade}[Atividade 3 — Apresentação dos Grupos e Validação Docente]

Com a planilha oficial projetada na tela para toda a sala:

\begin{enumerate}
\item Cada grupo terá \textbf{2 a 3 minutos} para apresentar a sua linha atualizada;
\item O grupo comentará:
  \begin{itemize}
      \item O estágio marcado na coluna \textbf{STATUS};
      \item Os avanços conquistados e a resolução das pendências apontadas no início da aula;
      \item A principal dificuldade enfrentada no momento.
  \end{itemize}
\item Os professores realizarão considerações pontuais e direcionamentos para o grupo.
\end{enumerate}

\end{boxAtividade}

% =====================================================
% PARTE 4
% =====================================================

\subsection{Parte 4 — Tarefa Prática 3: Definição da Sprint 1 de 2026-2 (20 minutos)}

\begin{boxAtividade}[Atividade 4 — Planejamento de Metas Imediatas]

Com base no estágio de maturidade indicado no \textbf{STATUS}, o grupo deverá cadastrar na coluna \textbf{PRÓXIMOS PASSOS} as \textbf{3 tarefas técnicas prioritárias} para as próximas 2 a 3 semanas:

\begin{itemize}
\item \textbf{Grupos em Diagnóstico ou Planejamento (ex.: JifTV, RasPorta):} Focar em recuperar ambiente/código legado, fechar escopo viável e rodar o protótipo mínimo;
\item \textbf{Grupos em Execução, Experimentos ou Consolidação (ex.: D20Hub, Scout de Vôlei):} Focar em refatoração de código, coleta de dados experimentais e escrita de Resultados no Overleaf.
\end{itemize}

\end{boxAtividade}

% =====================================================
% CHECKLIST
% =====================================================

\subsection{Checklist Final da Aula}

\begin{boxAtividade}[Verificação de Entregas da Aula]

Antes de sair do laboratório, certifique-se de que seu grupo concluiu as 4 ações:

\begin{itemize}
\item O link do \textbf{Pitch Vídeo (3 min)} foi cadastrado na planilha (para os 9 grupos pendentes);
\item O \textbf{Diagrama Canva} e a descrição do \textbf{Experimento 1} foram regularizados (\textit{JifTV} e \textit{RasPorta});
\item As colunas \textbf{AVANÇOS}, \textbf{DIFICULDADES} e a classificação padronizada de \textbf{STATUS} foram preenchidas;
\item O grupo apresentou sua linha em plenária e registrou os \textbf{PRÓXIMOS PASSOS} da Sprint 1.
\end{itemize}

\end{boxAtividade}

% =====================================================
% REFLEXÃO FINAL
% =====================================================

\subsection*{Reflexão Técnica}

A evolução de um projeto de engenharia exige transparência, acompanhamento contínuo e execução disciplinada. A planilha atualizada é o seu mapa de bordo para garantir um segundo semestre produtivo e de alto nível técnico!
